\section{Materials and Methods}
\label{sec:methods}

\begin{figure}[!t]
\centering
% Color definitions
\definecolor{mainborder}{rgb}{0.18, 0.24, 0.35}     % Sleek steel navy
\definecolor{mainfill}{rgb}{0.94, 0.95, 0.97}       % Elegant light slate gray
\definecolor{secborder}{rgb}{0.35, 0.43, 0.53}     % Soft slate/steel blue
\definecolor{secfill}{rgb}{0.98, 0.98, 0.99}       % Premium off-white/very light gray

\resizebox{\linewidth}{!}{%
\begin{tikzpicture}[
  >={Stealth[length=4pt]},
  mainbox/.style={
    rectangle, rounded corners=3pt, draw=mainborder, fill=mainfill,
    minimum height=1cm, minimum width=2.4cm,
    font=\small\bfseries, align=center, line width=0.8pt
  },
  detailbox/.style={
    rectangle, rounded corners=2pt, draw=secborder, fill=secfill,
    minimum height=0.7cm, minimum width=1.8cm,
    font=\footnotesize, align=center, line width=0.5pt
  },
  metricbox/.style={
    rectangle, rounded corners=2pt, draw=secborder, fill=secfill,
    minimum height=0.7cm, minimum width=2.8cm,
    font=\footnotesize, align=center, line width=0.5pt
  },
  arr/.style={->, thick, mainborder},
  link/.style={-, secborder, line width=0.5pt}
]

% ── Main pipeline (y = 0), compressed spacing for better scaling ──
\node[mainbox] (data)    at (0,   0) {Dataset\\[-1pt]{\footnotesize LLFF}};
\node[mainbox] (split)   at (3.2, 0) {Split\\[-1pt]Generation};
\node[mainbox] (train)   at (6.4, 0) {Training /\\[-1pt]Adaptation};
\node[mainbox] (eval)    at (9.6, 0) {Evaluation};
\node[mainbox] (out)     at (12.8,0) {Metrics};

\draw[arr] (data)  -- (split);
\draw[arr] (split) -- (train);
\draw[arr] (train) -- (eval);
\draw[arr] (eval)  -- (out);

% ── Split details (centered under split = 3.2) ──
\node[detailbox]        (unif) at (2.2, -1.7) {\textit{uniform}};
\node[detailbox]        (rand) at (4.2, -1.7) {\textit{random}};
\node[detailbox]        (seed) at (3.2, -2.7) {fixed seed (42)};
\draw[link] (unif.south) -- (seed.north -| unif.south);
\draw[link] (rand.south) -- (seed.north -| rand.south);
\draw[link] (split.south) -- (unif.north);
\draw[link] (split.south) -- (rand.north);

% ── Regime details (stacked under train = 6.4) ──
\node[detailbox] (zs)  at (6.4, -1.7) {Zero-Shot};
\node[detailbox] (tta) at (6.4, -2.5) {TTA {\scriptsize ($N\!\in\!\{3,6,10\}$)}};
\node[detailbox] (ps)  at (6.4, -3.3) {Per-Scene {\scriptsize (100\%)}};
\draw[link] (train.south) -- (zs);

% ── Eval detail (centered under eval = 9.6) ──
\node[detailbox] (loader) at (9.6, -1.7) {fixed\_indices};
\draw[link] (eval.south) -- (loader);

% ── Metric details (centered under out = 12.8) ──
\draw[link] (out.south) -- ++(0,-0.85);
\node[metricbox] (m_qual) at (12.8, -1.7) {PSNR, SSIM, LPIPS};
\node[metricbox] (m_eff)  at (12.8, -2.5) {FPS, VRAM};

% ── Environment annotation ──
\begin{scope}[on background layer]
  \node[draw=secborder!30, densely dashed, rounded corners=5pt,
        fill=black!1, inner sep=12pt,
        fit=(data)(split)(train)(eval)(out)(seed)(ps)(loader)(m_qual)(m_eff),
        label={[font=\scriptsize\itshape, text=secborder!70]below:
          Docker isolated GPU | NVIDIA GTX 1660 Super (6 GB) | CUDA 12.6 | $\sim$180h total execution}] {};
\end{scope}

\end{tikzpicture}%
}
% Caption and Label go AFTER the graphic/tikzpicture in IEEE format
\caption{Experimental execution protocol. The top boxes represent the pipeline stages; the bottom branches detail the configuration variants (split strategy, training regime) and the collected metrics. The entire environment was isolated via a Docker container with a reserved GPU.}
\label{fig:pipeline}
\end{figure}


This section presents the proposed methodology for systematically evaluating the trade-offs between computational performance and reconstruction fidelity of NeRF models under sparse-view conditions. Figure \ref{fig:pipeline} illustrates the proposed evaluation pipeline, which consists of three main stages: (i) dataset selection and sparse-view sampling; (ii) configuration and training of the selected neural rendering architectures according to their respective learning paradigms; and (iii) quantitative evaluation using standardized metrics to jointly characterize computational efficiency and reconstruction fidelity.

\subsection{Dataset \& Splits}
\label{sec:data}
The experiments were conducted on the LLFF (\textit{Local Light Field Fusion}) dataset, which comprises 8 real-world scenes (\textit{fern, flower, fortress, horns, leaves, orchids, room, trex}) captured with a forward-facing moving camera. Camera poses and depth bounds were obtained through COLMAP, following standard dataset procedures \cite{mildenhall2019locallightfieldfusion}. Furthermore, images were evaluated using specifically downscale of 8 resolution ($504 \times 378$).

The complete data splitting protocol follows the few-shot evaluation configuration established by Niemeyer et al. \cite{niemeyer2022regnerf}. Specifically, the testing split was generated by reserving images with indices that are multiples of 8 (\texttt{llffhold=8}), while the training views for the sparse regimes were strictly limited to $N \in \{3, 6, 10\}$. The selection of these training views was performed using two distinct strategies: (i) \textit{uniform}, where views are sampled via linear interpolation over the dataset indices, distributing them evenly along the camera trajectory (a strategy that does not generalize to other datasets, since LLFF indices intrinsically encode camera trajectory order); and (ii) \textit{random}, where views are randomly selected using a fixed seed (42) to ensure reproducibility. Comparing these two strategies allows us to evaluate the models' sensitivity to the spatial distribution of the input views, as demonstrated in Table \ref{tab:split}. Given the empirical advantage of uniform sampling (+0.54 dB PSNR on the \textit{fern} scene), this strategy was adopted for all subsequent experiments.

\begin{table}[!htbp]
  \centering
\caption{Impact of the training view selection strategy (\textit{uniform} versus\ \textit{random}) on average visual quality. Evaluated under the TTA regime with $N=10$ views on the \textit{fern} scene.}
\label{tab:split}
\resizebox{\columnwidth}{!}{%
\begin{tabular}{llcccc}
\toprule
\textbf{Model} & \textbf{Split} & \textbf{PSNR (dB) $\uparrow$} & \textbf{SSIM $\uparrow$} & \textbf{LPIPS $\downarrow$} & \textbf{$\Delta$PSNR} \\
\midrule
\multirow{2}{*}{PixelNeRF}
  & Uniform   & 6.98   & 0.047    & 0.846    & \multirow{2}{*}{+0.54} \\
  & Random    & 6.44   & 0.042    & 0.840    &                        \\
\midrule
\multirow{2}{*}{GNT}
  & Uniform   & 10.34  & 0.179    & 0.738    & \multirow{2}{*}{+0.15} \\
  & Random    & 10.19  & 0.184    & 0.742    &                        \\
\midrule
\multirow{2}{*}{Instant-NGP}
  & Uniform   & 13.57  & 0.231    & 0.627    & \multirow{2}{*}{+0.06} \\
  & Random    & 13.51  & 0.229    & 0.630    &                        \\
\midrule
\multirow{2}{*}{Nerfacto}
  & Uniform   & 12.04  & 0.227    & 0.809    & \multirow{2}{*}{+0.11} \\
  & Random    & 11.93  & 0.226    & 0.865    &                        \\
\bottomrule
\end{tabular}%
}
\end{table}

Beyond the core optimization hyperparameters, specific pipeline configurations were strictly enforced to guarantee stability and computational efficiency across the few-shot regimes. Notably, camera pose refinement was disabled for Nerfacto to prevent the catastrophic geometric collapse frequently observed in sparse-view scenarios. To minimize memory footprint and accelerate the Test-Time Adaptation phase, mixed-precision training was enabled across all architectures. For the hash-based models (Instant-NGP and Nerfacto), the multiscale hash table capacity was constrained to a maximum size of $2^{16}$, with backgrounds explicitly set to white. Furthermore, structural hyperparameters for GNT were locked to 48 coarse and importance samples per ray, a transformer depth of 8, and a network width of 64, ensuring reproducible structural bounds during volume rendering and following the default configuration from the original paper \cite{t2023attentionnerfneeds}.

\subsection{Models \& Training Regimes}
The proposed strategy evaluates representative models for novel view synthesis under few-shot constraints, encompassing distinct architectural paradigms. The generalizable approaches selected were: (i) {PixelNeRF \cite{yu2021pixelnerfneuralradiancefields}, which utilizes Convolutional Neural Networks (CNNs) to extract spatial features and inject them into the renderer via projection; and (ii) Generalizable NeRF Transformer (GNT) \cite{t2023attentionnerfneeds}, which replaces both the feature extractor and the differentiable renderer with cross-attention modules. As comparative baselines, two highly efficient per-scene architectures were chosen: Instant-NGP \cite{muller2022}, which replaces traditional MLPs with a multiscale hash table encoding, and Nerfacto \cite{Tancik2023}, an aggregate model combining state-of-the-art heuristics for rapid convergence and high fidelity.

The experimental execution was structured across three distinct regimes: (i) Zero-Shot: The ideal paradigm for extreme compute constraints, performing instant inference using pre-trained weights without any local optimization on the novel scene; (ii) Test-Time Adaptation (TTA):} An intermediate regime where the pre-trained generalizable models are fine-tuned using the sparse input views ($N \in \{3, 6, 10\}$, as reported in Table \ref{tab:room_n_views}); (iii) Per-Scene (100\%): Representing the upper bound for visual fidelity, this regime optimizes the baseline models from scratch using the entirety of available images for a specific scene. It is crucial to note that during the Test-Time Adaptation (TTA) phase for the generalizable models, the core feature extraction networks (CNN and ViT) were kept frozen, updating only the terminal rendering layers. Preliminary ablation tests revealed that unfreezing the entire network under extreme data sparsity ($N \le 10$) led to catastrophic forgetting and severe overfitting to the training poses, actively degrading novel view synthesis performance.

\begin{table}[!htbp]
\centering
\caption{Impact of $N$ training views on PSNR (dB) on the \textit{room} scene.}
\label{tab:room_n_views}
\footnotesize
\begin{tabular}{lccc}
\toprule
Model & N=3 $\uparrow$ & N=6 $\uparrow$ & N=10 $\uparrow$ \\
\midrule
GNT (ZS) & 12.96 & 12.85 & 12.85 \\
GNT (TTA)       & 11.78 & 12.85 & 11.58 \\
PixelNeRF (ZS) & 7.35 & 7.15 & 7.29 \\
PixelNeRF (TTA) & 7.24 & 7.26 & 7.13 \\
\bottomrule
\end{tabular}
\end{table}

\begin{table}[!htbp]
\centering
\caption{PSNR (dB) under varying learning rates ($\eta$). Ablation conducted on the \textit{fern} scene under the TTA regime with $N = 10$ views.}
\label{tab:ablation_lr}
\footnotesize
\begin{tabular}{llc}
\toprule
Model (Regime) & Learning Rate ($\eta$) & PSNR (dB) $\uparrow$ \\
\midrule
PixelNeRF (TTA)         & $10^{-4}$ & 11.19 \\
PixelNeRF (TTA)         & $5\times10^{-4}$ & 12.23 \\
PixelNeRF (TTA)         & $10^{-3}$ & 9.42 \\
PixelNeRF (ZS)          & N/A         & 6.64 \\
\midrule
GNT (TTA)         & $10^{-4}$ & 12.96\\
GNT (TTA)         & $5\times10^{-4}$ & 12.96\\
GNT (TTA)         & $10^{-3}$ & 12.85\\
GNT (ZS)          & N/A         & 10.73 \\
\bottomrule
\end{tabular}
\end{table}

To ensure a strictly equitable evaluation protocol, all models were unified under the \textit{Nerfstudio} framework \cite{Tancik2023}. Instant-NGP and Nerfacto are natively supported by the framework. Conversely,  custom integration wrappers for PixelNeRF and GNT were developed and publicly released, manually porting their default components into semanthic DataManager, Pipeline, and Model components to comply with Nerfstudio's modular interface without altering their native inference logic. This unified environment leverages \textit{tiny-cuda-nn} \cite{tinycudann} for GPU acceleration and standardizes the logging infrastructure. 

For consistency across the evaluation, all models were optimized using Adam with a standard learning rate of $\eta = 10^{-3}$, zero weight decay, 512 rays per batch in training and 128 in the evaluation step. Optimization budgets were scaled according to the regime: TTA configurations ran for 2500 iterations, while per-scene optimizations ran for 30000 iterations to ensure full convergence. Although per-architecture tuning of $\eta$ would yield individually higher peak metrics (as evidenced by PixelNeRF's 2.81 dB improvement under $\eta = 5 \times 10^{-4}$, Table \ref{tab:ablation_lr}), such tuning introduces a confounding variable that conflates hyperparameter sensitivity with architectural capability. To isolate architecture as the primary experimental variable, a uniform $\eta = 10^{-3}$ was enforced across all models. This represents a deliberate methodological constraint: absolute per-model performance is sacrificed in favor of a controlled comparison baseline. This criterion factors into the subsequent TTA stability analysis. 


The complete wrapper integration code is publicly available at \url{https://github.com/XXXX/XXXX/} %\url{https://github.com/Felipe-gsilva/Few-Shot-NeRFs/}
. Additionally, the modular wrappers can be downloaded directly as packages from PyPI and piwheels \url{https://www.piwheels.org/project/XXXX/} (%\url{https://www.piwheels.org/project/nerfstudio-gnt/} 
and \url{https://www.piwheels.org/project/XXXX/}
%\url{https://www.piwheels.org/project/nerfstudio-pixelnerf/}
), while the reproducible experiment scripts and environment configurations are hosted at \url{https://github.com/XXXX/XXXX/}
%\url{https://github.com/Felipe-gsilva/nerf-ann-paper/}
.

\subsection{Evaluation \& Metrics}
\label{sec:metrics}

The evaluation framework assesses three complementary dimensions: visual fidelity, computational efficiency, and memory consumption. Final metrics were computed exclusively on the held-out test set (section \ref{sec:data}), ensuring no overlap with the training views. 

To quantify visual quality, this project used three standard novel view synthesis metrics: (i) PSNR (Peak Signal-to-Noise Ratio), to measure pixel-wise fidelity; (ii) SSIM (Structural Similarity Index Measure) \cite{wang2004image}, to evaluate structural coherence; and (iii) LPIPS (Learned Perceptual Image Patch Similarity) \cite{zhang2018unreasonable}. Following the protocol established by Niemeyer et al. \cite{niemeyer2022regnerf}, LPIPS was computed using Visual Geometry Group (VGG) features to accurately quantify perceptual distance.

Operational efficiency was measured through the inference rendering rate in frames per second (fps) across the test cameras, as did Kulhanek et al. \cite{kulhanek2025nerfbaselinesconsistentreproducibleevaluation}. For memory evaluation,  both the static parameter footprint of the models and the dynamic VRAM allocation peak logged by PyTorch during runtime were analyzed.

To guarantee environmental isolation and reproducibility, all experiments were executed within Nvidia Docker containers, preventing external processes from adding noise to the performance metrics. The hardware setup consisted of an NVIDIA GeForce GTX 1660 Super (6GB VRAM) running CUDA 12.6.